\section{The long-wave check and a map back to the paper}
\label{sec:reading}
\subsection{Why Whitham theory appears near the end}
The paper uses a second approximation method as a consistency check.
Whitham modulation theory describes a locally periodic nonlinear wave
whose parameters change on scales much longer than its period and
wavelength. NLS assumes small amplitude and a narrow band. Whitham's
approach need not start with a small-amplitude carrier, but its modulation
must be very long. These are overlapping, rather than identical, assumptions.

For a longitudinal perturbation with $\lambda>0$ tending to zero,
\eqref{eq:consistent-root} gives
\begin{equation}\label{eq:long-wave}
 \Omega\simeq\frac\lambda2+\frac32\lambda A_0^2
              \mathbin{\pm}i\frac{\lambda A_0}{2\sqrt2}.
\end{equation}
The $\lambda^3$ linear-dispersion correction vanishes in this limit, and
$\lambda^2/K=|\lambda|$ tends to zero. The real correction proportional
to $\lambda A_0^2$ survives. This is the stated long-wave result (3.14).
For negative $\lambda$, the growing branch uses $|\lambda|$ in its
positive growth rate.

\subsection{Averaged Lagrangian: the minimum needed for (3.15)}
A Lagrangian is a function whose variation produces evolution equations.
For a nearly periodic wave, average it over the rapid phase to obtain
$\mathcal{L}(\omega,k)$ after the carrier-amplitude relation has been
incorporated. You do not need a full course in variational mechanics to
understand the particular check Dysthe makes.

Let the slowly varying phase be $S(x,t)$, so $k=S_x$ and $\omega=-S_t$.
Equality of mixed derivatives implies
\[
 k_t+\omega_x=0.
\]
Varying the averaged action with respect to $S$ yields
\[
 \partial_t\mathcal{L}_{\omega}-\partial_x\mathcal{L}_k=0.
\]
Linearize both equations about constant $\omega,k$. For a modulation
$e^{i(\lambda x-\Omega t)}$, the first gives
$\delta\omega=c\,\delta k$ with $c=\Omega/\lambda$.
Substitute into the second to get
\begin{equation}\label{eq:whitham}
 c^2\mathcal{L}_{\omega\omega}
 +2c\mathcal{L}_{\omega k}+\mathcal{L}_{kk}=0.
\end{equation}
This is (3.15), in the variables appropriate to that calculation.
The second derivatives are entries of the Hessian of $\mathcal{L}$.
If $\mathcal{L}_{\omega\omega}\ne0$, the roots are
\[
 c=\frac{-\mathcal{L}_{\omega k}\pm
 \sqrt{\mathcal{L}_{\omega k}^2-
       \mathcal{L}_{\omega\omega}\mathcal{L}_{kk}}}
 {\mathcal{L}_{\omega\omega}}.
\]
Negative discriminant means complex modulation speeds and an exponentially
growing branch. In PDE language, the real characteristic speeds of the
averaged system cease to exist.

The approximate averaged Lagrangian supplied in (3.16) is
\[
 \mathcal{L}=\frac{\rho g}{2k^2}(\sigma^2-\sigma^3-\sigma^4),
 \qquad \sigma=\frac{\omega^2}{gk}-1.
\]
Here $\omega$ denotes the \emph{nonlinear} carrier frequency, with
$\omega^2=gk[1+(ka)^2+\cdots]$. Differentiate the function with respect
to $\omega,k$ before evaluating it on the carrier relation. If you first
set $\omega^2=gk$, you incorrectly set $\sigma=0$ and erase the nonlinear
information the check needs.

Dysthe reports that this calculation agrees with (3.14) at the retained
small-amplitude order. That is a cross-check in the common long-wave,
small-steepness limit. The polynomial (3.16) is itself an approximation;
it is not an exact nonlinear dispersion law. The paper's historical
discussion of a long-wave instability boundary near $ka=0.34$ concerns
the separate finite-amplitude Whitham calculation, not a threshold proved
by its fourth-order envelope expansion.

\subsection{Read the actual paper in three passes}
\begin{enumerate}
\item \textbf{Physical argument.} Read pp.~105--106 and the conclusion on
pp.~113--114. Explain the chain: modulation changes averaged momentum
flux; a mean current develops; the current shifts phase; the altered
phase response changes modulational growth.
\item \textbf{Structure of the model.} Read Section~2, identifying the
interior equation, boundary conditions, harmonic expansion, scaling, and
mean-flow closure. Use the table below. Treat the nonlinear-product
algebra as a coefficient-collection calculation rather than memorizing it.
\item \textbf{Derivation and evidence.} Work through Section~3 with
Section~\ref{sec:stability} open. Check the determinant, the truncation,
and each figure axis. Finish with the long-wave overlap check.
\end{enumerate}

\begin{longtable}{@{}p{.22\linewidth}p{.49\linewidth}p{.22\linewidth}@{}}
\toprule Paper location & What the step does & Background here\\\midrule
\endhead
Introduction, pp.~105--106 & Compares small-amplitude NLS and very-long-wave
Whitham regimes; motivates one more asymptotic order. & Sections~1, 5, 6, 10\\
(2.1)--(2.3), p.~106 & States interior and free-surface fluid equations.
& Section~3\\
(2.4)--(2.5), p.~107 & Taylor expands surface conditions at $z=0$.
& Section~3\\
(2.6)--(2.8), p.~107 & Separates carrier, harmonics, and mean; finds leading
dispersion and group transport. & Sections~2, 4, 5\\
(2.9)--(2.11), pp.~107--108 & Averages to obtain induced mean potential
and elevation. & Section~7\\
(2.12)--(2.13), p.~108 & Uses Laplace's equation to eliminate envelope
vertical derivatives. & Appendix~A.1\\
(2.14)--(2.16), p.~108 & Collects the fundamental harmonic of nonlinear
products and its mean-flow interaction. & Sections~2, 5\\
(2.17)--(2.18), p.~109 & Forms an intermediate envelope equation and
nondimensionalizes it. & Section~6; A.2\\
(2.19)--(2.21), p.~109 & Reduces time derivatives; states fourth-order
equation and mean-flow closure. Check the coefficient issue. & Sections~7--8\\
(3.1)--(3.5), p.~110 & Introduces real amplitude/phase perturbations and
linearizes their mean flow. & Sections~6, 9\\
(3.6)--(3.7), p.~110 & Eliminates the mean field and solves a mode
determinant to consistent order. & Section~9\\
(3.8)--(3.12), pp.~111--112 & Extracts longitudinal instability band,
maximum growth, and frequencies. & Section~9\\
Figures 1--3, pp.~111--112 & Compares different observables with
full-equation calculations and some measurements. & Section~9\\
(3.13), pp.~112--113 & Studies oblique modulation and the NLS hyperbola.
& Sections~6, 9\\
(3.14)--(3.16), p.~113 & Checks the very-long-wave limit using an averaged
Lagrangian. & Section~10\\
(4.1)--(4.3), pp.~113--114 & Isolates mean-current feedback as the key
growth correction. & Sections~7--8\\
\bottomrule
\end{longtable}

\subsection{What to retain as the paper's conclusion}
The paper gives an improved \emph{approximate} description of weakly
nonlinear deep-water gravity-wave modulation. Its central stability
mechanism is the wave-induced mean current. The asymptotic growth and
frequency predictions are supported by the paper's reported comparisons
for particular observables, and by an overlap check with long-wave theory.

It does not derive breaking dynamics, prove a uniform error bound for all
packets, establish a large-bandwidth model, or validate every later
equation called Dysthe. In particular, the original two-horizontal-direction
temporal potential-envelope formulation is distinct from a later
one-coordinate spatial free-surface-envelope formulation.

\checkpoint{Explain why a smaller growth-rate error in Figure~2 does
not prove accurate phase prediction or accurate breaking dynamics.}
\clearpage
